\section{CONCLUSION}\label{sec:conclusion}

This article studied the pre-capture approach to a target that tumbles fast enough for its docking port to face the chaser only once per spin period, a period comparable to the deadline of the task. The proposed framework assigns the long-horizon decisions to a learned policy and the constrained motion to an execution layer of planners and MPC, and connects the two layers in both directions. RWG expresses the decisions in terms of the sweep-cone geometry: the latitude of the line of sight is regulated onto the cone swept by the port, the policy sets the closing pace and the timing of the encounter, and entries are planned toward the predicted port motion and re-timed within the same pass. ECC returns to every decision the entry opportunities that the execution layer can realize, keeps, corrects, or replaces the proposed intent by one rule, and trains the policy on the coordinated closed loop with the proposal stored, so that the critics learn the consequences of proposals under coordination with standard SAC losses.

In six-degree-of-freedom simulations of time-limited pre-capture, ECC captured $85.4\%$ of the test tasks against $73.8\%$ for learning-guided MPC on the same execution chain, $63.9\%$ with TD3, and $44.3\%$ for MPC alone. The gain consists of tasks captured in time, at the same completion time and propellant per capture, and it is concentrated in the tighter deadlines, where few passes remain. RWG raised the captures of a fixed intent from $64$ to $103$ of $192$ tasks, and entry re-timing added a small positive gain. ECC kept its capture rate under a slower spin and under a stronger nutation outside the training distribution, where the margin over learning-guided MPC grew to $22.4$ and $14.4$ percentage points.

The results are obtained in simulation with the relative state measured without noise, and the execution layer predicts the target attitude at constant angular velocity with an isotropic inertia estimate. Each learning method was trained with three seeds, and the confidence intervals reflect the variation over initial conditions. Of the $192$ test initial conditions, $191$ were captured by at least one method or training run, whereas ECC captured $152$ in all three runs and missed $19$ in all three; the remaining gap lies mainly in the consistency of learning. Future work will combine the coordination with onboard estimation of the target motion under sensing noise and improve the consistency of the learned policies across training runs.
